\documentclass[a4paper,12pt]{article}
\usepackage[utf8]{inputenc}
\usepackage[russian]{babel}
\usepackage{amsmath}
\usepackage{amsfonts}
\usepackage{geometry}
\usepackage{hyperref}
\usepackage{booktabs}

\geometry{left=2.5cm, right=1.5cm, top=2cm, bottom=2cm}

\title{\textbf{ПРОЕКТНОЕ ПРЕДЛОЖЕНИЕ}\\
\large Автоматизированная установка высокоинтенсивного конвективного теплообмена для циклического охлаждения жидких сред в замкнутом объеме}
\author{
    \textbf{Команда:} 3 человека \\\\
    \begin{tabular}{ll}
        1. Митюхин Иван, Б01-505 & \texttt{@Dercxh} \\
        2. Муругов Андрей, Б01-505 & \texttt{@murugov2007} \\
        3. Фафурин Егор, Б01-505 & \texttt{@Salmon258} \\
    \end{tabular} \\\\
    \textbf{Ментор:} Тавитов Азамат Георгиевич
}
\date{\the\year~г.}

\begin{document}

\maketitle

\section{Актуальность и идея проекта}
В бытовых условиях и сфере общественного питания часто возникает задача экстренного снижения температуры бутилированных напитков и жидких сред за минимальное время. Традиционные методы (статическое охлаждение в морозильных камерах) обладают низкой эффективностью из-за высокого термического сопротивления пограничного слоя жидкости у стенок тары. 

Проект направлен на создание компактного устройства, использующего контролируемое вращение тары по продольной оси в ледосолевой среде (метод \textit{Spin Chilling}). За счет создания вынужденной конвекции и развитого турбулентного режима внутри емкости время охлаждения сокращается с 60 минут до нескольких минут.

\section{Цель проекта}
Создать рабочий прототип мобильного экспресс-охладителя баночной тары объемом до 0.5~л, способного за счет автоматизированного гидродинамического воздействия существенно снижать температуру жидкого продукта с +25°C до +4°C в экспресс-режиме.

\section{Основной функционал}
\begin{itemize}
    \item \textbf{Регулируемый силовой привод:} вращение цилиндрической тары по продольной оси со скоростью до 150~об/мин (для текущего прототипа) с возможностью ШИМ-управления;
    \item \textbf{ШИМ-стабилизация вращения:} поддержание стабильного числа Рейнольдса для непрерывного перемешивания внутренних слоев жидкости без дегазации;
    \item \textbf{Быстрозажимной адаптивный узел:} надежная осевая фиксация алюминиевых банок стандартных геометрических размеров на базе деталей, оптимизированных методом 3D-печати;
    \item \textbf{Влагозащищенное управление:} запуск цикла охлаждения по таймеру одной физической кнопкой с автоматическим отключением силового привода;
    \item \textbf{Эндотермический термостат:} поддержание отрицательной температуры охлаждающей среды (до --5°C) за счет оптимизации фазового перехода водно-солевого раствора со льдом.
\end{itemize}

\section{Конструкция прототипа}
\begin{itemize}
    \item Микроконтроллер \textbf{ESP32 DevKitC} (32-бит, встроенные интерфейсы Wi-Fi/Bluetooth);
    \item Двухканальный H-мостовой драйвер двигателя \textbf{TB6612FNG};
    \item Высокомоментный электродвигатель постоянного тока \textbf{JGA25-370} (12V, 150 RPM);
    \item Импульсный понижающий DC-DC преобразователь \textbf{Mini560 (выход 5V)} для безопасного питания логики микроконтроллера;
    \item Герметичный цифровой датчик температуры \textbf{DS18B20} (контактный, в защитной капсуле);
    \item Тактовая кнопка запуска цикла;
    \item Быстрозажимной патрон (3D-печать);
    \item Герметичный охлаждающий бак (термоизолированный);
    \item Сетевой блок питания постоянного тока \textbf{12V} (не менее 1.5--2А).
\end{itemize}

\section{Принцип работы}
В охлаждающий бак заливается ледосолевой раствор с температурой до --5°C. Емкость с жидкостью фиксируется в быстрозажимном узле, после чего по нажатию кнопки запускается вращение. Вынужденная конвекция внутри тары разрушает термический пограничный слой, тепло эффективно отводится через стенки в охлаждающую среду. 

Микроконтроллер ESP32 через драйвер TB6612FNG управляет длительностью цикла, осуществляет плавный пуск/остановку двигателя по ШИМ, опрашивает цифровой датчик температуры DS18B20 для мониторинга среды и автоматически отключает привод по истечении 60 секунд.

\section{Задачи проекта}
\begin{itemize}
    \item Разработка механизма быстрой фиксации банок разного диаметра на валу;
    \item Сборка и разводка электрической схемы на базе платы расширения для ESP32 и драйвера TB6612FNG;
    \item Изготовление и герметизация охлаждающей ванны;
    \item Программирование микроконтроллера ESP32 (в среде PlatformIO/Zed на C++):
    \begin{enumerate}
        \item Обработка сигнала аппаратной кнопки пуска в режиме энергосбережения;
        \item Управление двигателем по ШИМ с использованием периферии LEDC в ESP32;
        \item Опрос шины OneWire датчика температуры DS18B20;
        \item Реализация неблокирующего таймера автоматического отключения;
    \end{enumerate}
    \item Тестирование охлаждения (замер температуры до/после при базовых 150 об/мин);
    \item Экспериментальный анализ эффективности и сравнение с естественным охлаждением в морозильной камере.
\end{itemize}

\section{Элементная база}
\begin{itemize}
    \item Плата разработки \textbf{ESP32 DevKitC};
    \item Мотор-редуктор постоянного тока \textbf{JGA25-370 12V 150RPM};
    \item Драйвер шаговых/DC двигателей \textbf{TB6612FNG};
    \item Стабилизатор \textbf{Mini560 (5V)};
    \item Цифровой термометр \textbf{DS18B20};
    \item Импульсный адаптер питания \textbf{12V};
    \item Термоизолированный контейнер и льдосолевая смесь;
    \item Тактовая кнопка, резистор 4.7 кОм, соединительные провода, плата расширения (\textit{breakout board});
    \item 3D-печатные детали (патрон, крепления).
\end{itemize}

\section{Ожидаемый результат}
Стабильно функционирующий сетевой прототип установки на базе 32-битной платформы, позволяющий проводить циклическое охлаждение баночной тары 0.5~л. Экспериментальное подтверждение эффективности метода \textit{Spin Chilling} при текущей конфигурации оборотов мотора, фиксация динамики падения температуры.

\section{Распределение ролей}
\begin{itemize}
    \item \textbf{Митюхин Иван} и \textbf{Фафурин Егор} --- проектирование механики, 3D-печать адаптивных узлов, гидроизоляция и сборка корпуса.
    \item \textbf{Муругов Андрей} --- проектирование схемотехники, разводка цепей питания и логики, написание встроенного ПО на C++ (PlatformIO/Zed).
\end{itemize}

\section{Этапы выполнения}
\begin{enumerate}
    \item Проектирование и разводка схемы на макетной плате (Текущий этап);
    \item Написание базовой прошивки (ШИМ, прерывания кнопки, OneWire);
    \item 3D-печать зажимного патрона и проектирование узла вывода вала;
    \item Сборка финального стенда в герметичном баке;
    \item Серия тестовых испытаний, замеры эффективности и финальная демонстрация.
\end{enumerate}

\section{Анализ существующих аналогов}
Существующие бытовые охладители (например, термоэлектрические кулеры) работают медленно (15–30 минут). Промышленные ротационные охладители (\textit{spin chillers}) используются в HoReCa, но стоят от сотен тысяч рублей и не предназначены для мобильного индивидуального использования. 

Предлагаемое устройство значительно доступнее аналогов за счет использования современной недорогой электроники (ESP32), применения аддитивных технологий (3D-печать) и отсутствия массивных компрессорных систем.

\end{document}